% Part: applied-modal-logics
% Chapter: epistemic-logic
% Section: relational-models

\documentclass[../../../include/open-logic-section]{subfiles}

\begin{document}

\olfileid{aml}{el}{rel}

\olsection{关系模型}

认知逻辑的基本语义概念与通常的模态逻辑相同。
关系模型仍然由一个可能世界的集合和一个赋值构成，该赋值决定了哪些命题变元在哪些世界上为“真”。
如果我们只处理单一主体，那么和通常一样，我们只有一个可达关系。
然而，如果我们处理的是多主体认知逻辑，那么单一的可达关系就变成了一组可达关系，主体符号集~$G$ 中的每个~$a$ 对应一个可达关系。

一个\emph{关系模型}由一个可能世界的集合和赋值组成，世界之间通过二元可达关系——每个主体对应一个——相联系，而赋值则决定了哪些命题变元在哪些世界上为真。

\begin{defn}
  多主体认知语言的\emph{模型}是一个三元组
  $\mModel{M} = \tuple{W, R, V}$，其中
  \begin{enumerate}
  \item $W$ 是一个非空的“世界”集合，
  \item 对每个 $a \in G$，${R}_a$ 是 $W$ 上的一个二元可达关系，且
  \item $V$ 是一个函数，为每个命题变元~$p$ 指派一个可能世界的集合 $V(p)$。
  \end{enumerate}
  若 $R_a ww'$ 成立，我们说 $w'$ \emph{从 $w$ 经由 $a$ 可达}。若 $w \in V(p)$，我们说 $p$ \emph{在 $w$ 上为真}。
\end{defn}

其机制与正规模态逻辑的机制相同，只是增加了更多的可达关系。
对于给定的主体，我们通常将其可达关系解释为刻画其信息状态的某些方面。
例如，我们常将 $R_a ww'$ 视为表示 $w'$ 与~$a$ 在~$w$ 处的信息相一致。
换言之，在~$w$ 处，主体~$a$ 无法区分世界 $w$ 和世界~$w'$。

\end{document}